% BHU PYQ -- Auto-generated & error-corrected
\documentclass[11pt,a4paper]{article}

\usepackage[a4paper,top=2.5cm,bottom=2.5cm,left=2.8cm,right=2.8cm,headheight=14pt]{geometry}
\usepackage{amsmath,amssymb}
\usepackage[T1]{fontenc}
\usepackage[utf8]{inputenc}
\usepackage{lmodern}
\usepackage{microtype}
\usepackage{fancyhdr}
\usepackage{enumitem}
\usepackage{hyperref}
\usepackage{lastpage}
\usepackage{parskip}

\hypersetup{colorlinks=true,urlcolor=black,linkcolor=black,
  pdftitle={STA/STB-604: Reliability -- BHU PYQ},pdfauthor={Banaras Hindu University}}

\pagestyle{fancy}\fancyhf{}
\renewcommand{\headrulewidth}{0.4pt}\renewcommand{\footrulewidth}{0.4pt}
\fancyhead[L]{\small   \textit{Banaras Hindu University}}
\fancyhead[R]{\small   \textit{STA/STB-604: Reliability}}
\fancyfoot[C]{\small Page  \thepage\ of \pageref{LastPage}}


\newlist{parts}{enumerate}{2}
\setlist[parts,1]{label=(\alph*),leftmargin=*,itemsep=5pt,topsep=3pt}
\setlist[parts,2]{label=(\roman*),leftmargin=*,itemsep=3pt,topsep=2pt}

\newcommand{\pts}[1]{\hfill   \textit{[#1]}}


\begin{document}


\begin{center}
  {\large\bfseries BANARAS HINDU UNIVERSITY}\\[2pt]
  {\normalsize B.A./B.Sc. (Hons.) Semester VI Examination, 2013-14}\\[6pt]
  \rule{0.8\linewidth}{0.5pt}\\[6pt]
  {\large\bfseries Paper No. STA/STB-604: Reliability}\\[4pt]
  {\normalsize Time: 3 Hours\hfill Maximum Marks: 70}
\end{center}
\rule{\linewidth}{0.4pt}
\smallskip



\noindent   \textit{Instructions: Answer any   \textbf{five} questions. All questions carry equal marks. Symbols have their usual meanings.}
\bigskip

\medskip


\noindent   \textbf{Question 1.}

\begin{parts}
  \item Explain the various censoring schemes (Type-I, Type-II, and random censoring) and obtain the corresponding likelihood function for each, with examples. \pts{7}
  \item Write short notes on:
  
\begin{parts}
    \item Conditional probability of failure
    \item System configuration (series, parallel, $k$-out-of-$m$)
    \item Hazard rate
  \end{parts}\pts{7}
\end{parts}

\medskip\rule{\linewidth}{0.2pt}

\medskip


\noindent   \textbf{Question 2.}

\begin{parts}
  \item What do you mean by aging of a system? Discuss various failure-time distributions based on the aging effect (IFR, DFR, IFRA) with examples. \pts{7}
  \item Obtain the maximum likelihood estimate of the survival function and hazard rate in the case of Type-I censoring, where the probability density function is given by:
  \[
    f(x; \theta) = \frac{1}{ \theta} e^{-x/ \theta}, \quad x > 0, \quad  \theta > 0.
  \]\pts{7}
\end{parts}

\medskip\rule{\linewidth}{0.2pt}

\medskip


\noindent   \textbf{Question 3.}

\begin{parts}
  \item Show that for a coherent system, $h(\mathbf{p})$ is strictly increasing in each component reliability $p_i$. Also explain that for an $n$-component series system, the reliability importance of the most reliable component is the least, and for a parallel system the opposite holds. \pts{7}
  \item Obtain the UMVUE of the survival function for $f(x; \theta) =  \theta e^{- \theta x}$, $x \geq 0$, $ \theta > 0$, in the case of a complete sample. \pts{7}
\end{parts}

\medskip\rule{\linewidth}{0.2pt}

\medskip


\noindent   \textbf{Question 4.}

\begin{parts}
  \item For the exponential distribution $F(x; \theta) = 1 - e^{-x/ \theta}$, $x > 0$, $ \theta > 0$, show that under Type-II censoring with $r$ ordered observations $X_{(1)} \leq X_{(2)} \leq \cdots \leq X_{(r)}$, the statistic:
  \[
    T = \sum_{i=1}^{r} X_{(i)} + (n-r)X_{(r)}
  \]
  is the complete and sufficient statistic for $ \theta$. Hence obtain the UMVUE of $ \theta$. \pts{7}
  \item Prove the following results:
  
\begin{parts}
    \item $R(t) = \exp\!\left(-\int_0^t h(x)\,dx\right)$, where $h(x)$ denotes the hazard rate.
    \item For a series system of $n$ independent components: $R_{ \text{sys}}(t) = \prod_{i=1}^n R_i(t)$.
  \end{parts}\pts{7}
\end{parts}

\medskip\rule{\linewidth}{0.2pt}

\medskip


\noindent   \textbf{Question 5.}

Obtain the reliability function for a parallel structure of $n$ independent components where $r_i(t)$ denotes the reliability of the $i$-th component. Hence compute the reliability of a parallel system when components are independently and identically exponentially distributed with parameter $\lambda$. \pts{14}

\medskip\rule{\linewidth}{0.2pt}

\medskip


\noindent   \textbf{Question 6.}

\begin{parts}
  \item What do you mean by the structure function of a system? Explain the different structures (series, parallel, $k$-out-of-$m$, bridge) that can be constructed using $n$ components. \pts{7}
  \item Explain the terms coherent structure and irrelevant component. Show that for a coherent system of $n$ components: $h(\mathbf{p}) = \sum_{i=1}^n p_i(1-p_i)$.  \pts{7}
\end{parts}

\medskip\rule{\linewidth}{0.2pt}

\medskip


\noindent   \textbf{Question 7.}

\begin{parts}
  \item State the lack-of-memory (memoryless) property and prove it for the exponential distribution. \pts{7}
  \item What do you understand by paths and cuts of a system? Hence define minimal path sets and minimal cut sets with examples. \pts{7}
\end{parts}

\medskip\rule{\linewidth}{0.2pt}

\medskip


\noindent   \textbf{Question 8.}

\begin{parts}
  \item Discuss the Weibull lifetime model along with its uses and properties. \pts{7}
  \item Consider the exponential distribution $f(x; \theta) = \frac{1}{ \theta}e^{-x/ \theta}$, $x > 0$, $ \theta > 0$, under a Type-II failure-censoring scheme. Obtain the best critical region for testing:
  
\begin{parts}
    \item $H_0: R(t) = R_0(t)$ vs.\ $H_1: R(t) = R_1(t)$, where $R_0(t) > R_1(t)$.
    \item $H_0: h(t) = h_0(t)$ vs.\ $H_1: h(t) = h_1(t)$, where $h_0(t) > h_1(t)$.
  \end{parts}
  Compute the power function and comment on its nature. \pts{7}
\end{parts}

\vfill
\rule{\linewidth}{0.4pt}

\begin{center}\small
  \href{https://bhu-exam.vercel.app/}{Click here to visit our website}\\[4pt]
  \href{https://me-aryan.vercel.app/}{Click here to visit our contributor}\\[4pt]
  \href{https://quashan-me.vercel.app/}{Click here to visit our contributor}
\end{center}

\end{document}